% Source fingerprint: 707844febb841f9cc4f51b07d6b7a079b052745c6b3db1946fc16c58ddf39776
% T16: raw TeX cells preserved; no numeric highlighting.
\begin{table}[htbp]
\centering\small\renewcommand{\arraystretch}{1.18}\setlength{\tabcolsep}{4pt}
\caption[紧性判据在哪类空间成立]{紧性判据在哪类空间成立。每个等价关系均带上空间类别。闭子空间继承完备性与任意子空间继承相对拓扑是不同命题。}
\label{b1:tab:visual-t16}
\begin{tabularx}{\linewidth}{@{}>{\raggedright\arraybackslash}p{3.1cm}>{\raggedright\arraybackslash}p{4.7cm}>{\raggedright\arraybackslash}X@{}}
\toprule
空间类别 & 成立的关系 & 不能省去的限定 \\
\midrule
任意拓扑空间 & 紧性用开覆盖定义；连续像保持紧性 & 序列收敛未必能描述闭包，紧性不自动给出序列子列性质 \\
度量空间 & 紧、序列紧、完备且全有界相互等价 & 全有界不能仅改为有界 \\
完备度量空间 & 全有界子集的闭包紧 & 子集本身若不闭仍可能不紧 \\
度量空间的紧子集 & 闭且全有界 & 一般拓扑空间中由紧推出闭还需 Hausdorff 条件 \\
可分度量空间 & 具有可数拓扑基 & 一般可分拓扑空间不据此自动得到第二可数性 \\
\bottomrule
\end{tabularx}
\end{table}
